\section{Literature Review}
\label{sec:literature}

Automated sustainability analysis covers several tasks. These tasks are finding relevant text segments, identifying commitments, checking whether a statement is specific, retrieving evidence, and organising results. Each task gives a useful signal, but a language signal alone does not prove that a company has misled a reader. Misleading a reader about environmental performance is commonly called \emph{greenwashing}. Calamai et al. review these tasks and point out continuing problems with available datasets and evaluation methods \cite{calamai2025}. For this reason, the review below treats each task separately from the claim of verified greenwashing.

\subsection{Language models and training labels}

BERT is a language model that learns the meaning of a word from the words around the word \cite{devlin2019}. This general learning is called \emph{pretraining}. After pretraining, BERT can be trained further on labelled examples to classify text. Gururangan et al. show that extra training on text from the same field and task improves several classification tasks \cite{gururangan2020}. This finding supports adapting a model to sustainability reports. However, a lower language-model loss, which is the error measured during training, does not by itself show better principle classification.

Sentence-BERT turns each sentence into a \emph{vector}, which is a list of numbers that represents the meaning of the sentence \cite{reimers2019}. Two vectors can then be compared with \emph{cosine similarity}, a measure of how close two meanings are. Similarity helps find related wording, but similarity does not confirm that a statement is supported by evidence.

Manual labelling is slow, so \emph{weak supervision} uses rules to create labels automatically. Snorkel combines several labelling rules and estimates how uncertain the resulting labels are \cite{ratner2017}. GRASP uses a simpler method. Each label is assigned from the tags and concept names stored in the report. A wrong assignment can therefore be repeated across many dataset rows. Kapoor and Narayanan explain how \emph{data leakage}, where information from the test data reaches the training data, and weak evaluation design can make machine-learning results look stronger than the evidence supports \cite{kapoor2023}. This risk motivates four checks in GRASP: company overlap, repeated text, label sources, and saved-model loading.

\subsection{Climate topics, commitments, and communication quality}

ClimaText provides labelled sentences for detecting climate-related topics \cite{varini2020}. A topic label helps select relevant text, but the label does not show whether a claim is supported. ClimateBERT adds pretraining on climate text and is tested on climate-language tasks \cite{webersinke2021}. Bingler et al. use ClimateBERT to study selective and unclear reporting under the Task Force on Climate-related Financial Disclosures (TCFD) framework \cite{bingler2022}. A later study links \emph{cheap talk}, meaning climate statements with little substance, to climate initiatives, emissions growth, and reputation risk \cite{bingler2024}. These studies already connect the language of companies with performance information. The contribution of GRASP is therefore defined only for the BRSR workflow.

ClimateBERT-NetZero studies net-zero and reduction commitments using approximately 3,500 expert-labelled examples, target classification, and question answering \cite{schimanski2023}. Finding a target and extracting the details of the target are separate tasks. Vinella et al. study greenwashing risk with generated labels and a fine-tuned ClimateBERT model \cite{vinella2024}. Generated labels increase the training data, but the evaluation still depends on how the labels were produced. The same concern applies to the AI-assisted concept mapping in GRASP.

\subsection{Report-level tools and Indian disclosure studies}

ChatReport combines language-model help, text retrieved from reports, and expert feedback to support disclosure review \cite{ni2023}. ESGReveal prepares documents, asks standard questions, and uses language-model agents to extract structured information from ESG reports \cite{zou2025}. Both studies show the value of a connected review tool. GRASP starts from structured XBRL facts instead. XBRL is a standard digital format in which every value already carries a report identifier and numerical context. For these fields, extracting text from page layouts is not required.

Shankar and Xu describe a BERT classifier trained with expert labels to classify sentences in Indian sustainability reports \cite{shankar2025}. The source is an SSRN working paper, which is a paper shared before peer review. The paper was posted in 2025, and the manuscript is dated 2024. Garg et al. review BRSR disclosures from 1,012 Indian businesses for 2022--23 \cite{garg2025}. This institutional report gives Indian reporting context and a broad review of disclosure practices. The report is not a machine-learning detection benchmark.

\subsection{Design choices for the BRSR workflow}

The literature supports three design choices. Evidence should be connected to report identifiers, the way training labels are obtained should be recorded, and the calculation of a company score should be shown. GRASP applies these choices through a report database, a nine-principle classifier, review flags, and sector comparisons. The evaluation compares label sources, principle categories, and company groups. Table~\ref{tab:lit} lists selected modelling and disclosure studies in year order and states the design lesson from each study.

{\small
\begin{longtable}{@{}L{0.85cm}L{2.35cm}L{2.65cm}L{2.65cm}L{2.65cm}L{2.65cm}@{}}
\caption{Related research and lessons for GRASP. Each study uses a different task and dataset, so reported findings are not scores on one shared test.}\label{tab:lit}\\
\toprule
\textbf{Year} & \textbf{Objective} & \textbf{Methodology} & \textbf{Main findings} & \textbf{Scope / limitation} & \textbf{Inference drawn}\\
\midrule\endfirsthead
\multicolumn{6}{l}{\textit{Table \thetable\ continued}}\\
\toprule
\textbf{Year} & \textbf{Objective} & \textbf{Methodology} & \textbf{Main findings} & \textbf{Scope / limitation} & \textbf{Inference drawn}\\
\midrule\endhead
\midrule\multicolumn{6}{r}{\textit{Continued on next page}}\\\endfoot
\bottomrule\endlastfoot
2019 & Learn general text features: BERT \cite{devlin2019}. & Pretraining with context on both sides of a word, followed by labelled training. & One pretrained model supports several language tasks. & General pretraining does not define BRSR categories. & Add a nine-principle output layer and test the BRSR task.\\
2019 & Compare sentence meanings \cite{reimers2019}. & Paired sentence encoders and vector similarity. & Supports efficient search for related sentences. & Similar wording does not prove factual support. & Use similarity as a review signal.\\
2020 & Adapt a model to specialised text \cite{gururangan2020}. & Further pretraining on domain and task text. & Improves several reported classification tasks. & Gains depend on the training text and target task. & Record the adaptation text separately from the labelled split.\\
2020 & Find climate-topic sentences: ClimaText \cite{varini2020}. & Dataset with sentence-level climate labels. & Provides data for testing topic detection. & A topic label does not assess a claim's credibility. & Separate topic detection from evidence checks.\\
2021 & Model climate language: ClimateBERT \cite{webersinke2021}. & Further pretraining on climate paragraphs. & Improves the reported climate-language tasks. & Climate tasks do not cover every social or governance principle. & Test adaptation on the intended task.\\
2022 & Study climate-risk reporting \cite{bingler2022}. & ClimateBERT analysis of TCFD-related language. & Finds vague claims and selective reporting. & TCFD categories differ from BRSR principles. & Keep labels specific to the reporting framework.\\
2023 & Find net-zero and reduction targets \cite{schimanski2023}. & Expert labels, ClimateBERT, and question answering. & Finds targets and related details. & Detecting a target does not verify delivery. & Separate commitments from measured outcomes.\\
2023 & Support report review: ChatReport \cite{ni2023}. & Text retrieval, model responses, and expert feedback. & Links report analysis to source text. & Retrieved text and answers still need evaluation. & Show the source text with each output.\\
2024 & Study risk using generated labels \cite{vinella2024}. & Generated labels and ClimateBERT fine-tuning. & Demonstrates an early language-model approach. & Results depend on the generated-label rules. & Identify AI labels separately from human labels.\\
2024 & Compare cheap talk with company outcomes \cite{bingler2024}. & Text-based indicator and outcome analysis. & Finds links with emissions growth and negative coverage. & A statistical link does not verify every claim. & Check the context of numerical evidence.\\
2025 & Extract ESG report data: ESGReveal \cite{zou2025}. & Document preparation, standard queries, retrieval, and model agents. & Supports structured extraction and disclosure analysis. & Results depend on documents and model choice. & Use available structured facts before page extraction.\\
2025 & Classify Indian report sentences \cite{shankar2025}. & Expert-labelled text and supervised BERT. & Reports sentence-level greenwashing detection. & Working paper with different data and labels. & Avoid direct accuracy comparisons across studies.\\
2025 & Review Indian BRSR disclosures \cite{garg2025}. & Institutional review of 1,012 businesses for 2022--23. & Describes BRSR reporting practices. & A disclosure study, not a detection test set. & Consider Indian reporting conditions in comparisons.\\
2025 & Review greenwashing text research \cite{calamai2025}. & Survey of tasks, data, and methods. & Identifies separate tasks and testing gaps. & No single language task establishes greenwashing. & State exactly what each output measures.\\
\end{longtable}}